% Vowel Ambiguity Patterns — Thematic Synthesis
% !TEX program = xelatex
% !TEX encoding = UTF-8 Unicode

\section{Vowel Ambiguity Patterns}
\label{sec:vowel-ambiguity}

Every colonial source in this study exhibits vowel underdetermination, but the
patterns differ systematically. Understanding what is systematic vs.\ random
in the vowel ambiguity is essential for designing conversion workflows. The
core problem is that each recorder used a five-letter Latin alphabet (a, e, i,
o, u) to represent an Algonquian vowel system of 8--10 phonemes, including
length distinctions and nasalization that the Latin alphabet was not designed
to encode.

\subsection{Per-Source Vowel Ambiguity}

Table~\ref{tab:vowel-ambiguity} shows how many Algonquian phonemes each
grapheme maps to in each source. The number of possible mappings per
grapheme ranges from 1 (for the most consistent digraphs) to 4 (for the
most ambiguous single letters).

\begin{table}[h]
\centering
\caption{Per-source vowel ambiguity: number of Algonquian phonemes per grapheme}
\label{tab:vowel-ambiguity}
\small
\begin{tabular}{lcccccc}
\toprule
\textbf{Grapheme} & \textbf{Williams} & \textbf{Eliot} & \textbf{Wood} & \textbf{Winslow} & \textbf{Edwards} & \textbf{Fielding} \\
\midrule
\textit{a} & 3 (/ə, ɔː, aː/) & 2 (/a, ə/) & 2 (/a, ə/) & 2 (/a, ə/) & 2 (/a, aː/) & 1 (/a/) \\
\textit{e} & 4 (/iː, ɛ, ə, ∅/) & 3 (/iː, ɛ, ə/) & 2 (/ɛ, i/) & 2 (/ɛ, ə/) & 2 (/iː, i/) & 1 (/iː/) \\
\textit{i} & 3 (/ə, iː, ɪ/) & 2 (/iː, ɪ/) & 2 (/ɪ, i/) & 2 (/ɪ, i/) & 1 (/i/) & 1 (/i/) \\
\textit{o} & 4 (/ə, aː, uː, ɔː/) & 3 (/uː, ɔ, ə/) & 3 (/ɔ, o, ə/) & 2 (/ɔ, ə/) & 2 (/o, u/) & --- \\
\textit{u} & 2 (/ə, a/) & 2 (/ə, ʌ/) & 2 (/ə, ʌ/) & 1 (/ə/) & 2 (/ə, o/) & 1 (/u/) \\
\textit{ee} & 1 (/iː/) & 1 (/iː/) & 1 (/iː/) & 1 (/iː/) & 1 (/iː/) & --- \\
\textit{oo} & 1 (/uː/) & 1 (/uː/) & 1 (/uː/) & 1 (/uː/) & 1 (/uː/) & --- \\
\midrule
\textbf{Total phonemes} & 8--10 & 8--10 & 6--8 & 6--7 & 7--8 & 6 \\
\bottomrule
\end{tabular}
\end{table}

The table reveals a clear hierarchy of ambiguity. The most reliable
transcriptions are the English digraphs \textit{ee} and \textit{oo}, which
consistently map to /iː/ and /uː/ respectively across all sources. The most
ambiguous are Williams' \textit{e} (4 possible values) and \textit{o} (4
possible values), reflecting the fact that Narragansett had a larger vowel
inventory than the five Latin letters could represent.

\subsubsection{Williams' Vowel System in Detail}

Williams' orthography is the most ambiguous because Narragansett had the
largest vowel inventory among the SNEA languages documented. His simple
vowel mappings (from O'Brien's Chart B) illustrate the problem:

\begin{itemize}
  \item \textit{a} $\rightarrow$ /ə/, /ɔː/, /aː/ --- three-way ambiguity
    covering schwa, open-o, and long a
  \item \textit{ah} $\rightarrow$ /aː/, /ɔː/ --- long variants
  \item \textit{e} $\rightarrow$ /iː/, /ɛ/, /ə/, ∅ --- four-way ambiguity;
    silent word-final
  \item \textit{ê} $\rightarrow$ /iː/ --- circumflex = long i (consistent)
  \item \textit{i} $\rightarrow$ /ə/, /iː/, /ɪ/ --- three-way; often = schwa
    in unstressed
  \item \textit{o} $\rightarrow$ /ə/, /aː/, /uː/; after \textit{w}: /aː/,
    /ɔː/ --- broad range
  \item \textit{oo}, \textit{ô} $\rightarrow$ /uː/ --- consistent long u
  \item \textit{ie}, \textit{ee} $\rightarrow$ /iː/ --- consistent
  \item \textit{ei} $\rightarrow$ /iː/, /ə/ --- two-way
  \item \textit{ea} $\rightarrow$ /iː/, /ɛə/, /aː/ --- three-way
\end{itemize}

\subsubsection{Eliot's Vowel System}

Eliot's orthography is more systematic but still underdetermined. His five
vowel letters (a, e, i, o, u) must represent 8--10 Massachusett phonemes.
The key disambiguating feature is that Eliot sometimes uses vowel doubling
to mark length: \textit{ee} = /iː/, \textit{oo} = /uː/, \textit{aa} = /aː/.
However, this doubling is not applied consistently across the Bible
translation. Eliot's distinctive ⟨ꝏ⟩ (oo ligature, often mis-encoded as ∞)
represents long /uː/, contrasting with \textit{oo} which may represent a
different vowel quality or length.

\subsubsection{Wood's Vowel Instability}

Wood's vowel system is the least stable. His triple spelling for 'devil'
(\textit{abbomocho}/\textit{abomocho}/\textit{abamacho}) shows that the
same recorder, recording the same word, produced three different vowel
transcriptions. The stable consonants (\textit{m}, \textit{ch}) are reliable,
while the unstable vowels (\textit{o}/\textit{a}) indicate reduced or
unstressed vowels that English ears struggled to categorize. This pattern is
diagnostic: when a word appears with multiple spellings in Wood, the stable
consonants are reliable while the unstable vowels indicate schwa or
unstressed reduction.

\subsubsection{Edwards' Vowel System --- The /iː/ Correction}

Recent work (Goddard 2007--2008, confirmed by Masthay 2024) has established
that Edwards' \textit{e} frequently represents /iː/, NOT /eː/ or /ɛ/. This
is the most important correction to the interpretation of Edwards'
orthography. Evidence includes:

\begin{itemize}
  \item \textit{Muhhekaneew} --- the second vowel is /iː/, not /eː/
  \item \textit{Neesoh} 'two' --- cf.\ PA *niːšwi, Munsee /niːšə/
  \item \textit{Seepoo} 'river' --- cf.\ PA *siːpoːwi, Munsee /siːpəw/
  \item Goddard's 2010 phoneme inventory for Mahican lists NO /eː/ phoneme
\end{itemize}

The default mapping for Edwards' \textit{e} should be /i(ː)/ unless cognate
evidence specifically indicates a different value.

\subsection{Systematic vs.\ Random Variation}

A critical distinction for the conversion workflow is whether vowel variation
is systematic (reflecting real phonological or dialect differences) or random
(reflecting the recorder's inability to consistently categorize the same
sound).

\subsubsection{Random Variation: Wood's Triple Spelling}

Wood's \textit{abbomocho}/\textit{abomocho}/\textit{abamacho} for 'devil' is
the clearest example of random variation. The same recorder, in the same
publication, produces three different transcriptions of the same word. The
vowel in the second syllable alternates between \textit{o} and \textit{a},
suggesting a reduced vowel that Wood could not consistently categorize. This
type of variation can be averaged out: the stable reading across multiple
attestations is likely the correct vowel quality.

\subsubsection{Systematic Variation: Williams vs.\ Winslow}

The consistent difference between Williams' \textit{wétu} and Winslow's
\textit{witeo} for 'house' represents systematic variation that may reflect
real dialect differences. Williams records Narragansett (Rhode Island), while
Winslow records Pokanoket Wampanoag (Plymouth area). The \textit{e}/\textit{i}
alternation in the first syllable and the final \textit{u}/\textit{o}
alternation are consistent across multiple attestations, suggesting a real
phonological difference between the dialects rather than random transcription
error.

\subsubsection{Diagnostic Criteria}

To distinguish systematic from random variation, we apply the following
criteria:

\begin{enumerate}
  \item \textbf{Consistency across attestations:} If the same source
    consistently writes the same vowel for the same word across multiple
    occurrences, the variation is likely systematic. If the vowel varies
    unpredictably, it is likely random.

  \item \textbf{Cross-source agreement:} If two independent sources agree on
    a vowel value, the variation is likely systematic. If they disagree, the
    disagreement may reflect dialect differences (systematic) or recorder
    error (random).

  \item \textbf{PA correspondence:} If the PA reconstruction predicts a
    specific vowel reflex, the variation can be evaluated against the
    prediction. Variation that matches the prediction is systematic;
    variation that does not is likely random.

  \item \textbf{Phonological environment:} Variation in unstressed syllables
    is more likely to be random (schwa reduction). Variation in stressed
    syllables is more likely to be systematic.
\end{enumerate}

\subsection{Schwa Reduction Patterns}

Schwa (/ə/) is the most variably transcribed vowel across all sources.
Colonial English recorders had no consistent way to represent schwa, so it
appears as \textit{a}, \textit{e}, \textit{i}, \textit{o}, or \textit{u}
depending on context and recorder. The following patterns emerge:

\begin{itemize}
  \item \textbf{Williams:} Schwa is most often written as \textit{i} or
    \textit{u} in unstressed syllables. Word-final schwa is often written as
    \textit{e} (which may be silent).

  \item \textbf{Eliot:} Schwa is most often written as \textit{u} in
    unstressed syllables. Eliot's grammar describes \textit{u} as the
    ``obscure'' vowel, corresponding to schwa.

  \item \textbf{Wood:} Schwa alternates unpredictably between \textit{a},
    \textit{e}, \textit{i}, \textit{o}, and \textit{u}. The triple spelling
    for 'devil' shows \textit{o} and \textit{a} alternating for what is
    likely a schwa.

  \item \textbf{Winslow:} Schwa is most often written as \textit{e} in final
    position (e.g., \textit{ahhe}, \textit{witeo}).

  \item \textbf{Edwards:} Schwa is most often written as \textit{u} in
    unstressed syllables. The sequence \textit{uh} may represent /əh/ or
    breathy schwa in a checked syllable.

  \item \textbf{Prince-Speck:} Schwa is written as \textit{'} (apostrophe)
    in their phonetic key, but in the orthographic forms it appears as
    \textit{a}, \textit{e}, or \textit{u}.
\end{itemize}

The schwa reduction pattern is so consistent across sources that it can be
used as a diagnostic: when a vowel in an unstressed syllable varies across
attestations, the default interpretation should be schva.

\subsection{Nasalization Ambiguity}

Nasalization is one of the most poorly documented features in the colonial
sources. Proto-Algonquian *aː and *oː developed nasalized reflexes in SNEA
languages before certain consonants, but no colonial source marks
nasalization consistently.

\begin{itemize}
  \item \textbf{Williams:} Nasalization is implied but not marked. He uses
    nasal coda consonant digraphs (\textit{an}, \textit{aun}, \textit{om},
    \textit{on}) to imply nasal vowels, but these same sequences also
    represent oral /an/, /aun/, /om/, /on/. Disambiguating requires
    checking the Massachusett cognate (Eliot marks nasal /ô/ explicitly) or
    the PA reconstruction.

  \item \textbf{Eliot:} The only colonial source that marks nasalization
    explicitly, using \textit{ô}, \textit{û}, or an overline. However, this
    marking is not applied consistently across the Bible translation.

  \item \textbf{Wood, Winslow:} Nasalization is entirely unmarked.

  \item \textbf{Edwards:} Nasalization is not marked. PA *aː developed a
    nasal quality in Mahican (marked by Schmick with a tilde), but Edwards
    does not indicate this.

  \item \textbf{Prince-Speck:} Nasalization is not systematically marked.
    May occasionally indicate it through vowel quality choices (e.g.,
    \textit{o} for nasalized /ã/).

  \item \textbf{Fielding:} Modern Mohegan has lost nasal vowels, so the
    modern source provides no help for this feature.
\end{itemize}

The nasalization ambiguity is consequential because it is phonemic in SNEA
languages. In Narragansett, \textit{n now} (yes) and \textit{núnow} (no) are
distinguished only by nasality, and Williams collapses this distinction
orthographically in related forms. The conversion workflow must infer
nasalization from comparative evidence rather than reading it from the
orthography.
